\documentclass{article}
\usepackage{lmodern}
\usepackage{amsmath}
\usepackage{listings}
\usepackage{fullpage}
\usepackage{parskip}
\usepackage{xcolor}
\lstset{
	basicstyle=\ttfamily,
	columns=fullflexible,
	frame=single,
	breaklines=true,
	postbreak=\mbox{\textcolor{red}{$\hookrightarrow$}\space},
}
\begin{document}
\title{Experiment `p\_4\_2\_c\_cumulative\_totals\_array' Results}
\author{\today}
\date{}
\maketitle
\textbf{Experiment outcome: }SUCCESS
\\\textbf{Bad responses: }0
\\\textbf{Responses containing}~\texttt{assume}~\textbf{: }0
\\\textbf{Responses containing}~\texttt{expect}~\textbf{: }0
\\\textbf{Resolution attempts: }3
\\\textbf{Hard fails (resolution): }0
\\\textbf{Soft fails (resolution): }0
\\\textbf{Verification attempts: }3
\section*{Problem Specification}
\textbf{Problem name: }p\_4\_2\_c\_cumulative\_totals\_array
\\\textbf{Natural language statement: }null
\\\textbf{Method signature: }p\_4\_2\_c\_cumulative\_totals\_array(inputs: array<int>) returns (totals: array<int>)
\subsection*{Ensures}
\begin{itemize}
\item \texttt{totals.Length == inputs.Length}
\item \texttt{forall i :: 0 <= i < inputs.Length ==> totals[i] == (if i == 0 then inputs[0] else totals[i - 1] + inputs[i])}
\end{itemize}
\subsection*{Requires}
\begin{itemize}
\item \texttt{0 < inputs.Length}
\end{itemize}
\clearpage
\section*{GenAI interactions}
Below you will find all interactions between the `user' (program) and the `assistant' (OpenAI).
\subsection*{Program $\rightarrow$ GenAI}
\begin{lstlisting}
You are given the following task to perform in Dafny:

Write a method to satisfy the below requirements.

The signature should be:

p_4_2_c_cumulative_totals_array(inputs: array<int>) returns (totals: array<int>)

The method should respect the following contract:

requires 0 < inputs.Length, ensures totals.Length == inputs.Length, ensures forall i :: 0 <= i < inputs.Length ==> totals[i] == (if i == 0 then inputs[0] else totals[i - 1] + inputs[i])

Produce and show only the Dafny body of this method, including the curly braces that surround it. Do not show the signature nor contract. You must not use 'assume' nor 'expect' anywhere in your code.
\end{lstlisting}
\subsection*{GenAI $\rightarrow$ Program}
\textbf{System fingerprint: }null
\\\textbf{ID: }chatcmpl-DnLNBDqeH8HJyGG1tMvD9tay45fRM
\\\textbf{Model: }gpt-5.5-2026-04-23
\\\textbf{Created at: }1780652601
\\\textbf{Finish reason: }stop
\\\textbf{Usage: }348 tokens in, and 640 tokens out
\begin{lstlisting}
{
  totals := new int[inputs.Length];
  totals[0] := inputs[0];

  var i := 1;
  while i < inputs.Length
    invariant 1 <= i <= inputs.Length
    invariant totals.Length == inputs.Length
    invariant forall j :: 0 <= j < i ==> totals[j] == (if j == 0 then inputs[0] else totals[j - 1] + inputs[j])
  {
    totals[i] := totals[i - 1] + inputs[i];
    i := i + 1;
  }
}
\end{lstlisting}
\subsection*{Program $\rightarrow$ GenAI}
\begin{lstlisting}
When using dafny verify, the below error is emitted and verify fails:

problem_data/p_4_2_c_cumulative_totals_array/program.dfy(13,14): Error: this invariant could not be proved to be maintained by the loop
 Related message: loop invariant violation
   |
13 |     invariant forall j :: 0 <= j < i ==> totals[j] == (if j == 0 then inputs[0] else totals[j - 1] + inputs[j])
   |               ^^^^^^


Dafny program verifier finished with 1 verified, 1 error

Correct the error in code by altering only the method body. Produce and show only the Dafny body, including the curly braces that surround it. Do not show the signature nor contract. You must not use 'assume' nor 'expect' anywhere in your code.
\end{lstlisting}
\subsection*{GenAI $\rightarrow$ Program}
\textbf{System fingerprint: }null
\\\textbf{ID: }chatcmpl-DnLNbYzcdtpTyKgjCl27MsGjYUraQ
\\\textbf{Model: }gpt-5.5-2026-04-23
\\\textbf{Created at: }1780652627
\\\textbf{Finish reason: }stop
\\\textbf{Usage: }694 tokens in, and 3635 tokens out
\begin{lstlisting}
{
  totals := new int[inputs.Length];
  totals[0] := inputs[0];

  var i := 1;
  while i < inputs.Length
    invariant 1 <= i <= inputs.Length
    invariant totals.Length == inputs.Length
    invariant forall j :: 0 <= j < i ==> totals[j] == (if j == 0 then inputs[0] else totals[j - 1] + inputs[j])
  {
    var oldI := i;
    ghost var prev := totals[..oldI];

    assert forall j :: 0 <= j < oldI ==> prev[j] == (if j == 0 then inputs[0] else prev[j - 1] + inputs[j]) by {
      forall j | 0 <= j < oldI
        ensures prev[j] == (if j == 0 then inputs[0] else prev[j - 1] + inputs[j])
      {
        assert prev[j] == totals[j];
        if j == 0 {
          assert totals[j] == inputs[0];
        } else {
          assert 0 <= j - 1 < oldI;
          assert prev[j - 1] == totals[j - 1];
          assert totals[j] == totals[j - 1] + inputs[j];
        }
      }
    }

    totals[oldI] := totals[oldI - 1] + inputs[oldI];
    assert totals[..oldI] == prev;

    i := oldI + 1;

    assert forall j :: 0 <= j < i ==> totals[j] == (if j == 0 then inputs[0] else totals[j - 1] + inputs[j]) by {
      forall j | 0 <= j < i
        ensures totals[j] == (if j == 0 then inputs[0] else totals[j - 1] + inputs[j])
      {
        if j == oldI {
          assert j != 0;
          assert totals[j] == totals[j - 1] + inputs[j];
        } else {
          assert j < oldI;
          assert totals[j] == prev[j];
          if j == 0 {
            assert prev[j] == inputs[0];
          } else {
            assert 0 <= j - 1 < oldI;
            assert totals[j - 1] == prev[j - 1];
            assert prev[j] == prev[j - 1] + inputs[j];
          }
        }
      }
    }
  }
}
\end{lstlisting}
\subsection*{Program $\rightarrow$ GenAI}
\begin{lstlisting}
When using dafny verify, the below error is emitted and verify fails:

problem_data/p_4_2_c_cumulative_totals_array/program.dfy(24,10): Error: assertion might not hold
   |
24 |           assert totals[j] == inputs[0];
   |           ^^^^^^

problem_data/p_4_2_c_cumulative_totals_array/program.dfy(49,12): Error: assertion might not hold
   |
49 |             assert prev[j] == inputs[0];
   |             ^^^^^^


Dafny program verifier finished with 1 verified, 2 errors

Correct the error in code by altering only the method body. Produce and show only the Dafny body, including the curly braces that surround it. Do not show the signature nor contract. You must not use 'assume' nor 'expect' anywhere in your code.
\end{lstlisting}
\subsection*{GenAI $\rightarrow$ Program}
\textbf{System fingerprint: }null
\\\textbf{ID: }chatcmpl-DnLOpEtpB931e6Of4XDi18B7N0K8Y
\\\textbf{Model: }gpt-5.5-2026-04-23
\\\textbf{Created at: }1780652703
\\\textbf{Finish reason: }stop
\\\textbf{Usage: }1482 tokens in, and 2543 tokens out
\begin{lstlisting}
{
  totals := new int[inputs.Length];
  totals[0] := inputs[0];

  var i := 1;
  while i < inputs.Length
    invariant 1 <= i <= inputs.Length
    invariant totals.Length == inputs.Length
    invariant totals[0] == inputs[0]
    invariant forall j :: 1 <= j < i ==> totals[j] == totals[j - 1] + inputs[j]
  {
    var n := i;
    totals[n] := totals[n - 1] + inputs[n];

    assert n != 0;
    assert totals[0] == inputs[0];

    assert forall j :: 1 <= j < n + 1 ==> totals[j] == totals[j - 1] + inputs[j] by {
      forall j | 1 <= j < n + 1
        ensures totals[j] == totals[j - 1] + inputs[j]
      {
        if j == n {
          assert totals[j] == totals[j - 1] + inputs[j];
        } else {
          assert j < n;
          assert 1 <= j < i;
          assert totals[j] == totals[j - 1] + inputs[j];
        }
      }
    }

    i := n + 1;
  }

  assert i == inputs.Length;
  assert forall k :: 0 <= k < inputs.Length ==> totals[k] == (if k == 0 then inputs[0] else totals[k - 1] + inputs[k]) by {
    forall k | 0 <= k < inputs.Length
      ensures totals[k] == (if k == 0 then inputs[0] else totals[k - 1] + inputs[k])
    {
      if k == 0 {
        assert totals[k] == inputs[0];
      } else {
        assert 1 <= k < i;
        assert totals[k] == totals[k - 1] + inputs[k];
      }
    }
  }
}
\end{lstlisting}
\section*{Final Program}
\begin{lstlisting}
method p_4_2_c_cumulative_totals_array(inputs: array<int>) returns (totals: array<int>)
	requires 0 < inputs.Length
	ensures totals.Length == inputs.Length
	ensures forall i :: 0 <= i < inputs.Length ==> totals[i] == (if i == 0 then inputs[0] else totals[i - 1] + inputs[i])
{
  totals := new int[inputs.Length];
  totals[0] := inputs[0];

  var i := 1;
  while i < inputs.Length
    invariant 1 <= i <= inputs.Length
    invariant totals.Length == inputs.Length
    invariant totals[0] == inputs[0]
    invariant forall j :: 1 <= j < i ==> totals[j] == totals[j - 1] + inputs[j]
  {
    var n := i;
    totals[n] := totals[n - 1] + inputs[n];

    assert n != 0;
    assert totals[0] == inputs[0];

    assert forall j :: 1 <= j < n + 1 ==> totals[j] == totals[j - 1] + inputs[j] by {
      forall j | 1 <= j < n + 1
        ensures totals[j] == totals[j - 1] + inputs[j]
      {
        if j == n {
          assert totals[j] == totals[j - 1] + inputs[j];
        } else {
          assert j < n;
          assert 1 <= j < i;
          assert totals[j] == totals[j - 1] + inputs[j];
        }
      }
    }

    i := n + 1;
  }

  assert i == inputs.Length;
  assert forall k :: 0 <= k < inputs.Length ==> totals[k] == (if k == 0 then inputs[0] else totals[k - 1] + inputs[k]) by {
    forall k | 0 <= k < inputs.Length
      ensures totals[k] == (if k == 0 then inputs[0] else totals[k - 1] + inputs[k])
    {
      if k == 0 {
        assert totals[k] == inputs[0];
      } else {
        assert 1 <= k < i;
        assert totals[k] == totals[k - 1] + inputs[k];
      }
    }
  }
}
\end{lstlisting}
\section*{Total Token Usage}
\textbf{Input tokens: }2524
\\\textbf{Output tokens: }6818
\\\textbf{Reasoning tokens: }5544
\\\textbf{Sum of `total tokens': }9342
\section*{Experiment Timings}
\textbf{Overall Experiment} started at 1780652600974, ended at 1780652769159, lasting 168185ms (168.19 seconds)
\\\textbf{Iteration \#1} started at 1780652600975, ended at 1780652626693, lasting 25718ms (25.72 seconds)
\\\textbf{Iteration \#2} started at 1780652626693, ended at 1780652703361, lasting 76668ms (76.67 seconds)
\\\textbf{Iteration \#3} started at 1780652703362, ended at 1780652769159, lasting 65797ms (65.80 seconds)
\end{document}
